\section{A native forest boundary that preserves the banking question}
\label{exec13:grf-adapter}

A servicing analyst asks whether making a specified offer changes
twelve-month net receipts for eligible customers. The proposed treatment
is the offer, coded one, against no offer, coded zero. The outcome
includes the offer's cost and is measured in dollars per customer over
the stated horizon. A heterogeneous-effect forest can help estimate this
contrast under a declared identification model. Its C++ implementation
cannot decide whether the recorded offer was randomized, whether all
confounders were measured, or whether the analyst actually wants the
mean over customers, branches, or treated customers. Those choices must
precede engine dispatch and survive it.

The proposed first grf integration uses its maintained R interface in a
pinned local process, which in turn invokes the native C++ core. Direct
native embedding is a separate alternative. Inspection covers commit
\texttt{5bee99b51471}, dated 30 April 2026, whose package description
declares version 2.6.1. The snapshot therefore precedes the literature
cutoff; its October acquisition date is separately retained. The design
is not an installed adapter. The executed examples below are
project-owned arithmetic references; they neither train grf nor measure
its statistical performance
\cite{exec13_grf_runtime,exec13_grf_design,exec13_grf_oracle}.

\subsection{From an offer question to a local moment}

Let \(X\) contain information available before the offer,
\(W\in\{0,1\}\) denote assignment, and \(Y(w)\) denote the potential
twelve-month net receipt. With consistency, no interference, conditional
exchangeability and positivity, the target is \[
 \tau(x)=E[Y(1)-Y(0)\mid X=x],\qquad
 \theta=E[\tau(X)].
\] The outer expectation defines the customer population. An observed
treated subpopulation or an overlap-weighted population gives a
different target. The initial proposed profile accepts complete
observations from independent customers, a binary treatment, one fixed
outcome and a fixed horizon. It excludes repeated decision histories,
unmeasured confounding, interference, survey weighting and latent-graph
ambiguity that has not been resolved for this estimand. These are
restrictions of the proposed profile, not assertions that the package
has no other functionality.

Write \(m(x)=E[Y\mid X=x]\) and \(e(x)=E[W\mid X=x]\). The inspected R
wrapper either obtains nuisance predictions or uses supplied
predictions, then passes \(Y-\widehat m(X)\) and \(W-\widehat e(X)\) to
the native training boundary. A direct call that passes the original
vectors does not replicate that boundary. The causal binding selects the
instrumental trainer and sets the instrument coordinate to the treatment
coordinate. Consequently the inspected native prediction is a ratio of
locally weighted covariances, with the treatment variance in its
denominator \cite{exec13_grf_rpath,exec13_grf_native}.

For normalized nonnegative forest weights \(a_i(x)\), let
\(\widetilde Y_i\) and \(\widetilde W_i\) denote the residualized
values. The local calculation has the form \[
 \widehat\tau(x)=
 \frac{\sum_i a_i(x)(\widetilde Y_i-\overline Y_a)
                         (\widetilde W_i-\overline W_a)}
      {\sum_i a_i(x)(\widetilde W_i-\overline W_a)^2},
 \quad
 \overline Y_a=\sum_i a_i(x)\widetilde Y_i,
 \quad \overline W_a=\sum_i a_i(x)\widetilde W_i.
\] In the exact reference, equal weights on residual outcomes
\((-2,-1,1,2)\) and residual treatments \((-1/2,-1/2,1/2,1/2)\) give
slope three. Constant residual treatment gives zero denominator and the
reference refuses the calculation. Its numerical threshold is an
explicit arithmetic rule, not evidence that positivity holds in a bank
population. The native source divides by its local first-stage quantity;
the adapter must inspect finite outputs and retain weak-support
diagnostics instead of silently replacing a failed division with zero.

Honest tree splitting and out-of-bag prediction address particular reuse
of observations. They do not establish the causal premises or turn
customer rows from a common branch shock into independent samples.
Generalized forest theory also imposes regularity, subsampling, smooth
local identification and variance conditions; a named package and a
successful fit do not verify them \cite{grf2018}. The proposed first
profile retains the R nuisance path and its out-of-bag predictions.
Supplying an independently estimated nuisance vector requires a later
profile with explicit training identities, alignment and sampling-unit
split rules. A known randomized propensity may be supplied only with the
corresponding assignment declaration; it is never inferred from a
convenient constant.

\subsection{Why the average effect needs its own calculation}

The average of fitted conditional effects is not interchangeable with
the package's default augmented average-effect estimator. In the
inspected binary path, write \(\widehat\tau_i\) for the out-of-bag
conditional prediction. Its score is \[
 \Gamma_i=\widehat\tau_i+
 \frac{W_i-\widehat e_i}{\widehat e_i(1-\widehat e_i)}
 \{Y_i-\widehat m_i-\widehat\tau_i(W_i-\widehat e_i)\}.
\] This equals the familiar two-outcome-regression expression after
setting \(\widehat\mu_{0i}=\widehat m_i-\widehat e_i\widehat\tau_i\) and
\(\widehat\mu_{1i}=\widehat m_i+(1-\widehat e_i)\widehat\tau_i\). The
equality is algebraic. Consistency and uncertainty still require
appropriate nuisance estimation, overlap and sampling assumptions
\cite{exec13_grf_rpath}.

The reference uses four tuples
\((Y,W,\widehat m,\widehat e,\widehat\tau)\): \((0,0,1,1/4,2)\),
\((1,0,1,1/2,2)\), \((3,1,2,1/2,2)\), and \((5,1,2,3/4,2)\). Their
scores are \(8/3,0,2,16/3\). The mean score is \(5/2\), whereas the mean
conditional prediction is two. The C++ residual-score calculation and an
independent rational Python calculation using the two-regression formula
agree. This is an exact conversion check with deliberately supplied
numbers, not fitted evidence for double robustness or coverage
\cite{exec13_grf_oracle}.

For target weights \(q_i\), the corresponding estimator is
\(\widehat\theta=\sum_iq_i\Gamma_i/\sum_iq_i\). Choosing \(q_i=1\)
targets a customer-weighted population; choosing inverse branch sizes
gives each observed branch equal total weight. A branch with two
customer scores of ten dollars and a branch with four scores of forty
dollars produce a customer mean of thirty dollars and an equal-branch
mean of twenty-five dollars. A planning cohort of ten thousand customers
turns the five-dollar difference into fifty thousand dollars in a naive
projection. That calculation exposes the target change; it does not
justify applying an equal-branch mean to an arbitrary customer
portfolio.

The R interface documents that equal cluster weighting changes both
training sampling and average-effect weights, and disallows user sample
weights in that mode. Replacing the final mean alone cannot reproduce
training under that option. A native port would have to preserve both
mechanisms. The initial independent-customer profile therefore rejects
cluster and weight options; the independent-cluster extension is
specified for later qualification and must name whether customers or
clusters define its target \cite{exec13_grf_rpath}.

\subsection{Sampling uncertainty is not a numerical decoration}

For independent clusters indexed by \(g\), define
\(S_g=\sum_{i\in g}q_i(\Gamma_i-\widehat\theta)\), \(Q=\sum_iq_i\), and
let \(G\) count positive-weight clusters. The inspected all-sample
augmented average-effect path computes the variance expression \[
 \widehat V_{\rm cl}=\frac{G}{G-1}\frac{\sum_g S_g^2}{Q^2}.
\] This specifies a computation. Asymptotic validity needs an
appropriate independent-cluster experiment, sufficient effective
clusters and nuisance conditions; no fixed small cluster count
guarantees it. Dependence shared across branches through a calendar
factor lies outside that law.

For the six scores above, treating all rows as independent gives
variance forty dollars squared. Preserving two branch clusters while
retaining the customer target gives \(1600/9\) dollars squared. Giving
the two branches equal total weight gives variance 225 dollars squared
around the changed target of twenty-five dollars. The standard errors
are approximately 6.325, 13.333 and 15 dollars. Two clusters make a
useful arithmetic counterexample and supply no empirical
confidence-coverage qualification. A single positive-weight cluster is
refused. Zero total weight is separately refused
\cite{exec13_grf_oracle,exec13_grf_rpath}.

Conditional prediction variance is another quantity. The native
prediction strategy estimates it using grouped trees and a
variance-debiasing calculation; the binding copies the variance values
into \texttt{variance.estimates}. The R average-effect operation instead
returns \texttt{std.err}, already square-rooted. Squaring the wrong
field or taking its square root twice changes the units. The proposed
adapter keeps the raw field identity, method and units, then creates any
reported interval only under a declared inference profile. Pointwise
conditional intervals cannot become a simultaneous customer band or an
interval for their average through serialization. Likewise
\texttt{excess.error}, a finite-forest error diagnostic, is not a
substitute for sampling uncertainty
\cite{exec13_grf_native,exec13_grf_rpath}.

\subsection{Memory ownership, coordinates and model persistence}

The inspected \texttt{Data} class is a read-only view of a borrowed
contiguous column-major buffer. Element \((r,c)\) is read at offset
\(cn+r\), where \(n\) is the number of rows. Constructing the view from
a temporary vector does not transfer ownership. The host must keep the
allocation alive until all training or prediction threads have finished;
it must validate shape, index ranges and multiplication overflow before
the call. A three-row, two-column example has columns \((1,2,3)\) and
\((10,20,30)\). Correct layout reads twenty at zero-based coordinate
\((1,1)\); passing row-major bytes to the same index formula reads
three. Neither value is a memory error, so a mere successful call cannot
detect the scientific corruption
\cite{exec13_grf_memory,exec13_grf_oracle}.

Role coordinates must remain separate from predictor coordinates. The
native setters mark outcome, treatment, instrument and weight columns as
disallowed split variables. The host records the ordered covariate
dictionary, row identifiers, role indices and training/query
distinction. It cannot infer those roles from whichever columns happen
to be numeric. Missing values in covariates are supported by inspected
package paths, but are excluded by the initial profile so that an
undocumented imputation or missingness assumption is not introduced.
Conversion to double precision must preserve identifier strings
separately and refuse nonfinite numerical values rather than silently
dropping affected customers.

Model serialization also has ownership semantics. The inspected Rcpp
helper iterates by moving the unique pointer to each native tree into
its serialized representation. The header explicitly says the original
forest cannot be used after that operation. A direct embedding that
serializes and then predicts on the consumed forest is therefore
invalid. Either serialize an explicitly supported copy, or deserialize
the saved representation into a new object. Retain the engine commit,
ordered feature map, original training-data identity, nuisance vectors,
options and format version beside the forest. The selected source does
not establish a stable C++ ABI or a portable standalone model format.
These are admission questions for a native route, not problems solved by
choosing a file extension \cite{exec13_grf_memory}.

\subsection{A concrete request and result contract}

The operation is \texttt{estimate.effect}. Its generic parameter schema
accepts an engine identifier and estimator identifier, so the proposed
profile uses \texttt{grf-r-native-2.6.1} and
\texttt{binary-iid-aipw-all-v1} there. Detailed settings reside in an
immutable configuration artifact, carried through the existing artifact
list. That artifact binds the exact source commit, treatment contrast,
outcome/horizon, row and column roles, seed, thread count, tree options
and target. The adapter validates its own profile before dispatch;
generic JSON acceptance alone does not validate these scientific
restrictions \cite{exec13_grf_design}.

Library, CLI and MCP surfaces must resolve identical authorized data and
configuration identities. The context digest retains the estimand,
population, units, information clock, assumptions, sampling declaration
and graph. The full request identity additionally binds data bytes,
configuration, engine environment and numerical settings. Changing the
target population or turning branch weighting on produces a different
request; transport names cannot cause reuse of a previous result. Graph
endpoint marks and unresolved sets remain intact. A forest is not a
graph-identification service: if a retained model class gives different
effects and the query has no accepted identification argument, this
profile cannot silently select a graph.

Before runtime admission, discovery of this engine returns unavailable.
A recognized but unsupported sampling regime, continuous treatment,
dynamic policy request or unresolved graph requirement returns
unsupported with an actionable explanation. Nonfinite output, failed
subprocess execution or an invalid model artifact is a failure; no
finite effect accompanies it. A reviewed identification failure can
return not-identified with its evidence, but the adapter's inability to
analyze a model is not an impossibility proof. An overlap alarm produces
a refusal under the declared profile. Trimming or switching to an
overlap population requires a newly reviewed target and request.

After engineering admission, a computed scalar estimate keeps its dollar
unit, signed offer-minus-no-offer contrast and target identifier. A
diagnostic conditional-prediction artifact retains query-row coordinates
and explicitly denotes conditional means, not individual counterfactual
differences. The initial engineering result carries execution-only
qualification and no confidence claim. Conditional inference needs its
own reviewed uncertainty construction and evidence. A mathematically
unbounded or disconnected result from another method must be returned
with its topology; this forest profile does not approximate it by a
large finite interval.

\subsection{Deployment choices and the next informative comparison}

The package description requires R at least 3.5.0, C++17 and GNU make,
links Rcpp and RcppEigen, and imports additional packages including
Matrix, DiceKriging, lmtest and sandwich. These declarations are
minimum/source requirements, not a tested environment lock. The proposed
R process must record the installed R and dependency versions, source
and binary hashes, compiler and flags, thread settings and runtime
options. The inspected native CMake file collects both source and test
files into an executable; its presence does not supply a ready
production shared library or a reviewed installed ABI. Package metadata
declares GPL-3 and the inspected source headers state GPL version 3 or
later. Preserve file and dependency terms and review distribution
obligations before shipping; a subprocess boundary alone settles neither
licensing compatibility nor distribution duties
\cite{exec13_grf_runtime}.

The parent process bounds wall time, memory and output size and
publishes only a complete validated result atomically. Standard output
is the protocol channel; diagnostics have a bounded separate log. A
crash or timeout leaves a typed failure and its input identities. No
mid-tree recovery is promised. Replaying an immutable request starts
from its recorded seed and environment; bitwise equality across
different compilers or thread settings remains a measured property. One
worker per isolated R process is the conservative initial choice because
the inspected binding also writes global runtime context. Concurrency
inside a shared embedding requires its own review.

The proposal recommends the R/native route first because it retains the
inspected nuisance, score and aggregation paths. Direct native embedding
may later reduce interface overhead or fit a host's deployment
constraints, but would transfer those statistical and ownership
responsibilities to the host. A full algorithm port has the greatest
maintenance and conformance burden. Reject direct embedding for the
initial profile until it reproduces the R path; this is a scoped
engineering choice, not a claim that native integration is impossible.

The next comparison should freeze one environment and the complete
workload: data conversion, nuisance estimation, forest training,
average-effect scoring, serialization, output validation and failure
handling. Compare the maintained R route with a direct-native candidate
only under identical inputs, target, options and estimator formulas.
Exact fixtures first check coordinates, residualization, score
aggregation and refusal behavior. A separately registered generated-data
study can then examine coverage, availability, heterogeneity recovery
and resources, retaining weak-overlap and dependent cases in its
denominator. Pure arithmetic agreement cannot replace that study, and
local runtime cannot establish bank workflow savings. The useful
investment question is whether a native route improves the complete
workload enough to justify owning the additional statistical and
deployment machinery. \label{exec13:grf-adapter-end}

\section{A common interface preserves differences that matter}

The adapter comparison supports a common request and result vocabulary with deliberately different admitted operations. The Python route returns a symbolic functional under a supplied graph and explicit conditional-path restrictions. The Java route proposes a graph class under a specified discovery model and must retain failed tests before publication. The R/native route estimates a defined customer effect while preserving nuisance adjustment, aggregation and uncertainty conventions. A single effect field cannot stand in for these three services.

The local trials show why preservation extends beyond the estimator. A coordinate permutation, rounding at an integer boundary, a changed context, corrupt replay or incomplete publication can alter or conceal the scientific result. The observed refusals and successful recoveries establish particular behavior; the descriptor, input-access and error-classification findings identify repairs still needed for those versions. The next service design must carry both kinds of evidence forward. Its value is that a caller can tell what was computed, under which assumptions, what remains uncertain and which next action could resolve the limitation.
